\documentclass[11pt,a4paper]{article}
\input{../../sablony/pripravy/priprava.sty}
\usepackage{tikz,pgfplots}
\pgfplotsset{compat=1.18}
\nastavPripravu{1. ročník}{Kvadratické rovnice a nerovnice}{111}{Rovnice vedoucí na kvadratickou: substituce}
\begin{document}
\blok{Rovnice vedoucí na kvadratickou: substituce}
Bikvadratické rovnice a substituce $t=x^2$; podmínka $t\ge0$.
\begin{definiceBox}
\textbf{Bikvadratická rovnice:} $ax^4+bx^2+c=0$, $a\ne0$.
\end{definiceBox}
\blok{Metoda a návaznost na graf}
\begin{itemize}
\item Polož $t=x^2\ge0$, vyřeš kvadratickou rovnici pro $t$ a teprve potom se vrať k $x$.
\item Záporné hodnoty $t$ nepřinášejí žádné reálné $x$.
\end{itemize}
\blok{Společně řešené příklady}
\begin{prikladReseniBox}[title={Příklad 1}]
$x^{4} - 5x^{2} + 4=0$
\tcblower
\[
\begin{aligned}
t&=x^2\ge0\\[2pt]
0&=t^{2} - 5 \cdot t + 4\\[2pt]
t&\in\{1;4\}\\[2pt]
K&=\{-2;-1;1;2\}
\end{aligned}
\]
\end{prikladReseniBox}
\begin{prikladReseniBox}[title={Příklad 2}]
$x^{4} + 2x^{2} - 3=0$
\tcblower
\[
\begin{aligned}
t&=x^2\ge0\\[2pt]
0&=t^{2} + 2 \cdot t - 3\\[2pt]
t&\in\{1\}\\[2pt]
K&=\{-1;1\}
\end{aligned}
\]
\end{prikladReseniBox}
\begin{prikladReseniBox}[title={Příklad 3}]
$x^{4} + 3x^{2} + 2=0$
\tcblower
\[
\begin{aligned}
t&=x^2\ge0\\[2pt]
0&=t^{2} + 3 \cdot t + 2\\[2pt]
t&\in\varnothing\\[2pt]
K&=\varnothing
\end{aligned}
\]
\end{prikladReseniBox}
\blok{Samostatné procvičování}
\begin{zadaniReseniBox}[title={Úloha 1}]
$x^{4} - 10x^{2} + 9=0$
\tcblower
\[
\begin{aligned}
t&=x^2\ge0\\[2pt]
0&=t^{2} - 10 \cdot t + 9\\[2pt]
t&\in\{1;9\}\\[2pt]
K&=\{-3;-1;1;3\}
\end{aligned}
\]
\end{zadaniReseniBox}
\begin{zadaniReseniBox}[title={Úloha 2}]
$x^{4} - 4x^{2}=0$
\tcblower
\[
\begin{aligned}
t&=x^2\ge0\\[2pt]
0&=t^{2} - 4 \cdot t\\[2pt]
t&\in\{0;4\}\\[2pt]
K&=\{-2;0;2\}
\end{aligned}
\]
\end{zadaniReseniBox}
\blok{Volitelná rezerva}
\poznamka{Studenti si vymyslí obdobný příklad se změnou jednoho koeficientu a vysvětlí, jak se změní počet kořenů či řešení.}
\blok{Volitelná výzva -- dva kladné kořeny substituce}
\begin{zadaniReseniBox}[title={Rozšiřující příklad}]
$x^4-10x^2+9=0$.
\tcblower
\[\begin{aligned}
t&=x^2\ge0\\t^2-10t+9&=0\\t&=1\ \text{nebo}\ 9\\|x|&=1\ \text{nebo}\ |x|=3\\K&=\{-3;-1;1;3\}
\end{aligned}\]
\end{zadaniReseniBox}
\blok{Domácí úkol}
\begin{zadaniReseniBox}[title={Domácí příprava}]
\begin{enumerate}[label=\alph*),itemsep=2pt]
\item $x^{4} - 13x^{2} + 36=0$
\item $2x^{4} - 5x^{2} + 2=0$
\item $x^{4} - 2x^{2} - 8=0$
\item $x^{4} + 4x^{2} + 4=0$
\item $x^{4} - 16=0$
\item $x^{4} - x^{2} - 12=0$
\end{enumerate}
\tcblower
\begin{enumerate}[label=\alph*),itemsep=2pt]
\item
\[\begin{aligned}
t&=x^2\ge0\\[2pt]
0&=t^{2} - 13 \cdot t + 36\\[2pt]
t&\in\{4;9\}\\[2pt]
K&=\{-3;-2;2;3\}
\end{aligned}\]
\item
\[\begin{aligned}
t&=x^2\ge0\\[2pt]
0&=2 \cdot t^{2} - 5 \cdot t + 2\\[2pt]
t&\in\{\frac{1}{2};2\}\\[2pt]
K&=\{- \sqrt{2};- \frac{\sqrt{2}}{2};\frac{\sqrt{2}}{2};\sqrt{2}\}
\end{aligned}\]
\item
\[\begin{aligned}
t&=x^2\ge0\\[2pt]
0&=t^{2} - 2 \cdot t - 8\\[2pt]
t&\in\{4\}\\[2pt]
K&=\{-2;2\}
\end{aligned}\]
\item
\[\begin{aligned}
t&=x^2\ge0\\[2pt]
0&=t^{2} + 4 \cdot t + 4\\[2pt]
t&\in\varnothing\\[2pt]
K&=\varnothing
\end{aligned}\]
\item
\[\begin{aligned}
t&=x^2\ge0\\[2pt]
0&=t^{2} - 16\\[2pt]
t&\in\{4\}\\[2pt]
K&=\{-2;2\}
\end{aligned}\]
\item
\[\begin{aligned}
t&=x^2\ge0\\[2pt]
0&=t^{2} - t - 12\\[2pt]
t&\in\{4\}\\[2pt]
K&=\{-2;2\}
\end{aligned}\]
\end{enumerate}
\end{zadaniReseniBox}
\poznamka{Orientační čas: úvod 5 min, společné příklady 18 min, samostatně 15 min, kontrola a shrnutí 7 min. Rezerva je dobrovolná. Číselné parametry modelových aplikací nejsou skutečně naměřená data.}
\end{document}
